\section{Методы}
\label{sec:methods}

\subsection{Экспериментальная установка}
\label{sec:experimental_setup}
В ходе экспериментов были идентифицированы поведенчески селективные нейроны у мышей при свободном исследовании квадратной арены открытого поля, обогащённой разнообразными контекстными стимулами. Мыши многократно помещались в среду для исследования, а их нейронная активность регистрировалась с помощью кальциевого имиджинга in vivo через минископ.

\subsubsection{Животные}
\label{sec:animals}
В экспериментах использовались мыши линии C57BL/6 (оба пола, n=16) в возрасте 3--5 месяцев. До хирургического вмешательства животные содержались в стандартных лабораторных клетках (по 2--7 особей) при 12-часовом цикле свет/темнота со свободным доступом к корму и воде. После операции мыши содержались индивидуально, прочие условия содержания оставались без изменений. Все эксперименты проводились в светлую фазу суток (с 10:00 до 18:00).

\subsubsection{Хирургические процедуры и доставка вируса}
\label{sec:surgeries}
Все процедуры были одобрены Комиссией по биоэтике Московского государственного университета имени М.\,В. Ломоносова (заявка № 159-г, утверждена на заседании Комиссии по биоэтике № 154-д-р от 17.08.2023) и соответствовали приказу Министерства здравоохранения РФ № 267 и Руководству NIH по содержанию и использованию лабораторных животных.
Хирургические процедуры выполнялись в три этапа: инъекция вируса, имплантация GRIN-линзы и установка базовой пластины. Анестезия индуцировалась внутрибрюшинной инъекцией Золетила (40 мг/кг) и ксилазина (5 мг/кг). Дексаметазон (4 мг/кг) вводился подкожно за 5 минут до операции для уменьшения воспаления и предотвращения отёка мозга. Для предотвращения высыхания глаз применялся увлажняющий гель. Мыши фиксировались в стереотаксической рамке (Kopf, США), температура тела поддерживалась с помощью нагревательного коврика (Physitemp, США). Для вирусной инъекции 300 нл AAV1.CAG.GCaMP6s доставлялись унилатерально в дорсальную область CA1 (координаты: AP $-1.94$ мм, ML $+1.46$ мм, DV $-1.2$ мм, \cite{Paxinos2019}). Через неделю GRIN-линза диаметром 1 мм (GrinTech, Германия) имплантировалась на 200 мкм выше CA1. Линза крепилась к черепу с помощью цианоакрилатного клея (Henkel, Германия) и стоматологического цемента (Stoelting, США) и покрывалась Kwik-Sil (World Precision Instruments, США). На третьем этапе, через неделю после предыдущей операции, над местом имплантации устанавливалась алюминиевая базовая пластина для крепления минископа, фиксировавшаяся стоматологическим цементом. Для защиты линзы устанавливалась пластиковая крышка. После заключительной операции мыши возвращались в домашние клетки и получали три недели на восстановление.

\subsubsection{Поведенческая процедура}
\label{sec:behavioral_procedure}
Через две недели после последней операции мыши постепенно приучались к обращению экспериментатора и ношению минископа. В течение трёх последовательных дней каждую мышь помещали в домашнюю клетку с подсоединённым, но не записывающим минископом на 5 минут за сессию. Данный этап позволял животным адаптироваться к весу и присутствию устройства без дополнительного стресса. После привыкания животные тестировались в квадратной арене открытого поля (44$\times$44$\times$44 см, серый пластик; Open Science, Россия) в течение четырёх последовательных дней. Каждая сессия длилась 10 минут. Арена была спроектирована таким образом, чтобы обеспечить обогащённую контекстную среду, включая полупрозрачный пузырчатый силиконовый коврик на полу и характерное расположение четырёх напольных объектов: белого пенопластового конуса, флакона для клеточных культур, наполненного наполнителем для подстилки, синего деревянного цилиндра и жёлтого пластикового строительного блока в форме параллелепипеда с отверстием. На каждой стене арены размещался уникальный визуальный ориентир: (1) жёлтый круг с пятью маленькими синими точками, (2) синий треугольник, (3) чёрный квадрат с широкими белыми полосами и (4) составная фигура из синих трапеций с чёрным перевёрнутым треугольником и белым кругом в центре. На протяжении каждой пробы поведение животного записывалось с помощью видеокамеры, установленной сверху (Flir Chameleon3, KVR, Великобритания), а кальциевая активность CA1 одновременно регистрировалась через минископ. Видеопоток поведения и поток минископа синхронизировались в Bonsai (Bonsai Foundation, Великобритания) с помощью общих временных меток, что обеспечивало согласованные временные базы для анализа зависимостей с задержкой.

\subsubsection{Анализ поведения}
\label{sec:behavioral_analysis}
Трекинг и аннотация поведения выполнялись с помощью пакета Sphynx для анализа исследовательского поведения \cite{Plusnin1}. Для каждого видео позиции ключевых точек тела оценивались покадрово с использованием нейронной сети, обученной в DeepLabCut \cite{Mathis}. Координаты с оценкой достоверности ниже 0.95 исключались и восстанавливались методом кусочной кубической интерполяции по соседним кадрам с высокой достоверностью. Полученные траектории сглаживались фильтром Савицкого--Голея третьей степени с окном 0.25 секунды для центра тела и основания хвоста и 0.1 секунды для всех остальных точек. Арена была сегментирована на функционально различные пространственные зоны: пристеночную зону шириной 7 см вдоль периметра, центральную зону 30$\times$30 см, угловые зоны 7$\times$7 см и объектные зоны шириной 2.5 см вокруг каждого объекта. На основе движения животного и его положения относительно этих зон рассчитывались как непрерывные, так и дискретные поведенческие переменные. К непрерывным переменным относились:
декартовы координаты центра тела;
абсолютная скорость, рассчитанная по перемещению центра тела и сглаженная с окном 0.25 секунды;
направление тела (НТ), определяемое как угол от центра тела к центру головы;
направление головы (НГ), определяемое как угол от центра головы к кончику носа.
Дискретные поведенческие акты классифицировались на основе динамики движения и пространственных переходов. Локомоторные состояния определялись по текущей скорости: быстрая локомоция (>5 см/с), медленная локомоция (1--5 см/с) и покой (<1 см/с). Замирание определялось как одновременное снижение скорости центра тела (<1 см/с) и скорости кончика носа (<2 см/с). Стойка на задних лапах детектировалась по кратковременному уменьшению расстояния между задними конечностями и основанием хвоста. Вхождения в зоны и взаимодействия с объектами определялись по положению центра тела и кончика носа соответственно. Все дискретные переменные сглаживались медианным фильтром с окном 0.25 секунды для устранения шума и обеспечения минимальной продолжительности событий.


\subsubsection{Обработка данных кальциевого имиджинга}
\label{sec:calcium_processing}
Данные кальциевого имиджинга обрабатывались с помощью BEARMIND --- специализированного аналитического конвейера, разработанного авторами (см. раздел о доступности кода), построенного на основе фреймворка CaImAn
\cite{Pnevmatikakis2016}. Программное обеспечение было создано для поддержки сквозной обработки многосессионных записей минископа.
Каждая сессия начиналась с визуального осмотра для определения границ кадрирования.
Этапы предобработки включали коррекцию фона и удаление двигательных артефактов с помощью алгоритма NoRMCorre \cite{Pnevmatikakis2017}. Затем видеозаписи разлагались на пространственные компоненты и временные ряды активности методом ограниченной неотрицательной матричной факторизации (CNMF-E) \cite{Zhou2018}. Гиперпараметры CNMF-E подбирались с использованием заданных критериев качества (компактность пространственного компонента, отношение сигнал/шум (SNR) временного ряда и остаточная структура) и верифицировались визуально. Основные диапазоны параметров составляли:
\begin{itemize}
\item \texttt{gSig} (оценка размера нейрона): 3--6 пикселей
(при разрешении 3 мкм/пиксель),
\item \texttt{min-corr} (порог локальной корреляции):
0.85--0.95,
\item \texttt{min-pnr} (отношение пик/шум): 3--15.
\end{itemize}
Ключевые результаты нижележащих этапов оставались качественно стабильными в данном диапазоне параметров (анализ робастности), а все итоговые компоненты подвергались ручной проверке для исключения артефактов и дубликатов.
После разложения предполагаемые компоненты автоматически фильтровались с использованием порогов \texttt{rval-thr} = 0.95 и \texttt{min-SNR} = 3 в соответствии с
\href{https://caiman.readthedocs.io/en/latest/CaImAn\_Tips.html}{рекомендациями CaImAn}.
Все сохранённые компоненты визуально проверялись в BEARMIND для выявления и исключения артефактов и дубликатов. Только компоненты, прошедшие эту ручную проверку, классифицировались как нейроны.

Следы флуоресценции нормализовались как $dF/F$, а пространственные карты совмещались между сессиями с помощью CellReg \cite{Sheintuch2017} для продольного отслеживания индивидуальных нейронов.

\subsubsection{Детекция кальциевых событий}
\label{sec:calcium_event_detection}
Для извлечения дискретных кальциевых транзиентов из следов $dF/F$ использовался вейвлетный подход, ранее описанный в \cite{Neugornet2021}, с незначительными модификациями. Данный метод был выбран за его устойчивость к шуму и способность обнаруживать многомасштабные особенности в данных флуоресценции.

След $dF/F$ предварительно сглаживался свёрткой с гауссовым ядром ($\sigma =
8$), после чего преобразовывался с помощью непрерывного вейвлет-преобразования (НВП), что давало двумерный массив вейвлет-коэффициентов по времени и масштабу. Значимые события проявлялись как <<гребни>> локальных максимумов на смежных масштабах (\fig{fig:pc}A).

Использовались обобщённые вейвлеты Морса $\Psi_{\beta,\gamma}(\omega) =
\alpha_{\beta,\gamma}\omega^{\beta}e^{-\omega^{\gamma}}$ с параметрами $\gamma =
3$ и $\beta = 2$, соответствующими так называемым вейвлетам Эйри, которые обладают минимальной площадью Гейзенберга и хорошей временной локализацией.
Для каждого масштаба обнаруживались локальные максимумы амплитуды вейвлет-коэффициентов. Максимумы на наибольшем масштабе инициировали гребни, которые затем продолжались на более мелкие масштабы на основе временного перекрытия. При наличии нескольких кандидатов на данном шаге в пределах каждого гребня сохранялся только наибольший локальный максимум, а остальные инициировали новые гребни. Процесс продолжался вплоть до наименьшего масштаба, при этом гребни, не охватывающие все масштабы, отбрасывались.
После построения всех гребней применялся фильтр значимости, основанный на длительности гребня, амплитуде пика и масштабе, на котором достигался максимум. Гребни, прошедшие этот этап фильтрации, классифицировались как кальциевые события.

\subsection{Конвейер INTENSE}
\label{sec:intense_pipeline}

\subsubsection{Вычисление взаимной информации}
\label{sec:mi_computation}

Для оценки взаимной информации (ВИ) использовался метод взаимной информации на основе гауссовой копулы (GCMI)
\cite{gcmi, ma2011mutual}. Данный метод использует свойство, согласно которому взаимная информация зависит только от структуры зависимости (копулы) переменных, а не от их индивидуальных маргинальных распределений, что обеспечивает эффективное вычисление в замкнутой форме через ранговое преобразование. В отличие от гистограммных оценщиков, GCMI позволяет избежать проблемы выбора числа интервалов \cite{Ross2014}; в отличие от оценщиков на основе $k$ ближайших соседей \cite{Kraskov2004}, метод не требует большого объёма выборки. Он обеспечивает надёжную оценку ВИ как для пар непрерывных переменных, так и для пар <<дискретная--непрерывная>>.
Для двух случайных величин $X$ и $Y$ взаимная информация определяется как:
\begin{equation}
\label{eq:mi_def}
I(X;Y) = H(X) + H(Y) - H(X,Y)
\end{equation}
GCMI сначала преобразует переменные в пространство гауссовой копулы:
\begin{equation}
\label{eq:gcmi_transform}
\tilde{X} = \Phi^{-1}(F_X(X)), \quad \tilde{Y} = \Phi^{-1}(F_Y(Y))
\end{equation}
где $F_X$ и $F_Y$ --- эмпирические функции распределения, а
$\Phi^{-1}$ --- обратная функция стандартного нормального распределения.
\textbf{Случай <<непрерывная--непрерывная>>}
Для двух непрерывных переменных ВИ оценивается через корреляционную матрицу
$\mathbf{R}$ преобразованных переменных:
\begin{equation}
\label{eq:gcmi_continuous}
I_{GCMI}(X;Y) = -\frac{1}{2}\log_2|\mathbf{R}|
\end{equation}
\textbf{Случай <<дискретная--непрерывная>>}
Для дискретной переменной $X$ с $K$ различными значениями и непрерывной переменной $Y$
ВИ вычисляется через разложение энтропии:
\begin{equation}
\label{eq:gcmi_discrete}
I_{GCMI}(X;Y) = H(Y) - \sum_{k=1}^{K} p(X=k) \, H(Y | X=k)
\end{equation}
% matching code (gcmi.py:mi_model_gd).
где $H(Y)$ --- дифференциальная энтропия $Y$, оценённая методом GCMI, а каждая условная энтропия по классам $H(Y|X=k)$ оценивается по подмножеству данных, где
$X = k$, с использованием предположения о гауссовой копуле.
Для каждой пары <<нейрон--поведенческая переменная>> сила селективности количественно оценивалась как:
\begin{equation}
\label{eq:selectivity_strength}
S = I(X_{neural}; Y_{behavior})
\end{equation}


\subsubsection{Статистическая оценка и идентификация поведенчески селективных нейронов}
\label{sec:statistical_assessment}

Для оценки значимости взаимной информации (ВИ) между кальциевой активностью и поведением наблюдаемое значение сравнивалось с нулевым распределением, построенным на основе данных со случайным временным сдвигом \cite{skaggs1993information}.
А именно, нейронный сигнал $X(t)$ подвергался
циклическому сдвигу на случайную величину $\tau_i$, после чего ВИ пересчитывалась:

$MI_{\text{shuffle}}^{(i)} = I(X(t + \tau_i); Y(t))$

Данная процедура повторялась многократно (далее --- <<перемешивания>>)
для построения нулевого распределения значений ВИ, ожидаемых при отсутствии связи. Циклический временной сдвиг сохраняет структуру временной
автокорреляции, присущую как нейронным кальциевым сигналам, так и поведенческим переменным, при этом разрушая их временную связь \cite{Lancaster2018}.
Это принципиально важно, поскольку кальциевые индикаторы порождают сигналы с медленной динамикой (постоянные времени затухания 1--2 секунды для GCaMP6s) \cite{Chen2013}, а поведенческие переменные часто демонстрируют сильные автокорреляции (например, пространственное положение при непрерывной локомоции).
Простое случайное перемешивание разрушило бы эти автокорреляции, создавая нереалистично низкое нулевое распределение.


Полученное нулевое распределение моделировалось гамма-распределением с дополнительной массой в нуле (zero-inflated gamma distribution), которое явным образом учитывает вероятностную массу в нуле, характерную для нулевых распределений ВИ:

$P(MI = 0) = \pi, \quad P(MI = x \mid x > 0) = (1 - \pi) \cdot f_\Gamma(x; \alpha, \beta)$

где $\pi$ --- параметр дополнительной массы в нуле, оценённый как доля значений ВИ при перемешивании, равных или близких к нулю, а $\alpha$ (форма) и $\beta$ (масштаб) оценивались методом максимального правдоподобия только по ненулевым значениям ВИ при перемешивании. Гамма-компонента была выбрана, поскольку ВИ строго неотрицательна, а распределения ВИ при перемешивании нейронных данных устойчиво демонстрируют правостороннюю асимметрию.

Значимость наблюдаемой ВИ оценивалась с помощью функции выживания подобранной модели:

$p\text{-value} = (1 - \pi) \cdot [1 - F_\Gamma(MI_{\text{observed}}; \hat{\alpha}, \hat{\beta})]$

где $F_\Gamma$ --- функция распределения гамма-распределения. Масса в нуле не вносит вклада в хвостовую вероятность для положительных наблюдаемых значений. Данный параметрический метод обеспечивает более гладкую и надёжную оценку хвостовых вероятностей по сравнению с прямыми эмпирическими квантилями, что особенно важно при ограниченном числе перемешиваний.

Для идентификации пар <<нейрон--поведенческая переменная>> со статистически значимым информационным содержанием применялась двухэтапная процедура, обеспечивающая баланс между вычислительной эффективностью
и контролем ложноположительных результатов:

\textbf{Этап 1: Скрининг}

Для каждой пары <<нейрон--переменная>> генерировалось 100 случайных перемешиваний; сохранялись пары, удовлетворяющие условию:

$MI_{\text{observed}} > \max_{i=1}^{100} MI_{\text{shuffle}}^{(i)}$

Данный предварительный скрининг позволял исключить неинформативные пары на раннем этапе, снижая вычислительные затраты и число гипотез для поправки на множественные сравнения.

\textbf{Этап 2: Валидация}

Для сохранённых пар выполнялось 10\,000 дополнительных перемешиваний с применением двух критериев:

(а) \textbf{Непараметрический критерий}: наблюдаемая ВИ должна превышать 99.95-й процентиль распределения при перемешивании:

$MI_{\text{observed}} > Q_{0.9995}(MI_{\text{shuffle}})$

(б) \textbf{Параметрический критерий}: $p$-значение, полученное из подобранного гамма-распределения (описанного выше), должно быть ниже скорректированного порога значимости:

$p\text{-value} = (1 - \pi) \cdot [1 - F_\Gamma(MI_{\text{observed}}; \hat{\alpha}, \hat{\beta})]
< p_{\text{threshold}}$

Порог $p$-значения определялся методом Холма--Бонферрони для контроля семейной ошибки первого рода (FWER) на уровне $\alpha = 0.01$:

$p_{\text{threshold}}^{(k)} = \frac{\alpha}{m - k + 1}$

где $m$ --- число гипотез (пар, прошедших этап 1), а $k$ --- ранг упорядоченных $p$-значений.

Несмотря на статистическую значимость, некоторые пары <<нейрон--переменная>> демонстрировали слишком слабое временное согласование для содержательной интерпретации. Для исключения таких случаев применялся дополнительный фиксированный порог:

$MI_{\text{observed}} > MI_{\text{threshold}}$

где $MI_{\text{threshold}}$ был постоянным для всех сессий данного эксперимента. Для анализируемого набора данных гиппокампальных записей
$MI_{\text{threshold}} = 0.04$ бита.

Совместное использование непараметрического (рангового) и параметрического критериев обеспечивает робастность: ранговый критерий защищает от ошибок подгонки распределения, типичных при наличии выбросов в биологических записях, тогда как параметрическое $p$-значение обеспечивает большую статистическую мощность для обнаружения слабых, но устойчивых ассоциаций.
Данный консервативный подход снижает частоту ложноположительных результатов при разведочном анализе многомерных нейронных данных.

\paragraph{Ускорение пермутационного тестирования методом БПФ.}

Описанная выше двухэтапная процедура требует вычисления ВИ при тысячах циклических сдвигов для каждой пары <<нейрон--переменная>>. Наивная реализация пересчитывала бы ВИ независимо для каждого сдвига, что даёт стоимость $O(n_{\text{sh}}
\cdot n)$ на пару, где $n$ --- число временных кадров, а $n_{\text{sh}}$ --- число перемешиваний. Для типичного эксперимента с $N$ нейронами, $M$ поведенческими признаками,
$D$ шагами задержки и $n_{\text{sh}} = 10{,}000$ перемешиваниями такие вычисления становятся
практически невыполнимыми.

Используется тот факт, что циклические временные сдвиги соответствуют циклическим кросс-корреляциям, которые для \emph{всех} $n$ возможных сдвигов
могут быть вычислены одновременно через теорему о свёртке:
%
\begin{equation}
\label{eq:fft_xcorr}
C_{XY}[s] = \sum_{i=0}^{n-1} X[i] \, Y[(i+s) \bmod n]
           = \mathcal{F}^{-1}\!\bigl[\mathcal{F}[X] \cdot \mathcal{F}[Y]^{*}\bigr][s]
\end{equation}
%
где $\mathcal{F}$ обозначает дискретное преобразование Фурье, а $^{*}$ --- комплексное сопряжение. Это снижает стоимость с $O(n_{\text{sh}} \cdot n)$ до
$O(n \log n)$ для вычисления ВИ при всех сдвигах, после чего значение при любом конкретном сдвиге извлекается за $O(1)$ как элемент массива. Конкретная форма ВИ при каждом сдвиге зависит от типов переменных:

\textit{Случай <<непрерывная--непрерывная>>.}
Для копула-нормализованных гауссовых переменных ВИ сводится к
$I(X; Y_s) = -\tfrac{1}{2}\log_2(1 - r(s)^2)$, где $r(s)$ --- коэффициент корреляции Пирсона при сдвиге $s$. Корреляция для всех сдвигов вычисляется из
ур.~\ref{eq:fft_xcorr} как $r(s) = C_{XY}[s] / [(n-1)\,\sigma_X \sigma_Y]$,
что требует единственной пары вещественнозначных БПФ.

\textit{Случай <<непрерывная--дискретная>>.}
Для непрерывной переменной $Z$ и дискретной переменной $X$ с $K$ классами
условные суммы по классам при каждом сдвиге представляют собой циклические корреляции с индикаторными функциями:
$\sum_i Z[i] \cdot \mathbb{1}[X_{(i+s)} = k] = \mathcal{F}^{-1}[\mathcal{F}[Z]
\cdot \mathcal{F}[\mathbb{1}_{X=k}]^{*}][s]$.
Из этих сумм и аналогичных сумм $Z^2$ вычисляются условные дисперсии по классам и гауссовы энтропии для всех сдвигов одновременно, что даёт ВИ через
$I = H(Z) - \sum_k p(k)\,H(Z | X\!=\!k)$. Для этого требуется $2K+2$ БПФ
(для $Z$, $Z^2$ и $K$ индикаторных функций) вне зависимости от числа перемешиваний.

\textit{Случай <<дискретная--дискретная>>.}
Таблицы сопряжённости для всех сдвигов получаются через индикаторные кросс-корреляции:
$n_{ij}(s) = \mathcal{F}^{-1}[\mathcal{F}[\mathbb{1}_{X=i}] \cdot
\mathcal{F}[\mathbb{1}_{Y=j}]^{*}][s]$, что требует $K_X \cdot K_Y$ пар БПФ.

\textit{Многомерный случай.}
Для $d$-мерных поведенческих признаков (например, 2D-координаты положения) ковариационная матрица разделяется на инвариантные к сдвигу внутрипеременные блоки и зависящие от сдвига блоки кросс-ковариаций, вычисляемые с помощью $d$ БПФ. Затем ВИ вычисляется через замкнутые векторизованные формулы для определителей ($d \leq 3$).

INTENSE предвычисляет ВИ для всех $n$ сдвигов один раз для каждой пары <<нейрон--признак>> и сохраняет результат в кэше БПФ. Этот кэш повторно используется при оптимизации задержки
(раздел~\ref{sec:temporal_alignment}), скрининге на этапе~1 и
валидации на этапе~2 без повторных вычислений. Суммарная стоимость составляет $O(N \cdot M \cdot n \log n)$
для построения кэша плюс $O(N \cdot M \cdot n_{\text{sh}})$ для выборки сдвигов,
по сравнению с $O(N \cdot M \cdot n_{\text{sh}} \cdot n)$ без ускорения ---
что даёт ускорение в $100$--$2{,}500\times$ в зависимости от типов пар и делает
валидацию на этапе~2 с 10\,000 перемешиваниями вычислительно выполнимой для крупномасштабных
наборов данных. На практике полный двухэтапный конвейер для 500 нейронов и 20 признаков
при 10-минутной записи с частотой 20~кадров/с и 10{,}000 циклическими сдвигами завершается приблизительно за 15 секунд (Intel~i9, многоядерная параллелизация через joblib).

\subsubsection{Временное согласование и оптимизация задержки}
\label{sec:temporal_alignment}

Для вычисления оптимальных задержек между сигналами тестировался диапазон задержек от $-2$ до $+2$ секунд с шагом 0.05 секунды; выбиралась задержка, максимизирующая
взаимную информацию между кальциевым сигналом и поведенческой переменной. Все
последующие вычисления выполнялись аналогично описанному для случая нулевой задержки,
с тем отличием, что <<истинная>> взаимная информация определялась как значение, вычисленное при оптимальном временном сдвиге.

GCaMP6s характеризуется временем нарастания $\sim$100 мс и временем затухания 1--2 с, что создаёт систематические задержки между спайками и флуоресценцией \citep{Zhang2023}. Различные поведенческие переменные
могут иметь различные оптимальные задержки для одного и того же нейрона.

Окно поиска $\pm 2$ секунды обосновано: (1) максимальным горизонтом предсказания,
наблюдаемым в навигационных задачах ($\sim$2 секунды) \citep{Lian2022}; (2) полным затуханием
кальциевых транзиентов GCaMP ($\sim$2 секунды) \citep{Zhang2023}; (3) поведенческими временными масштабами
при принятии решений грызунами; (4) вычислительной выполнимостью при охвате релевантных
временных масштабов.

Для каждой пары <<нейрон--признак>> из сетки кандидатов выбиралась задержка, максимизирующая взаимную информацию;
затем статистическая значимость оценивалась при данной фиксированной
задержке с помощью пермутационного тестирования на основе циклического сдвига; поправка на множественные сравнения
применялась по парам <<нейрон--признак>>, а не по задержкам.


\subsubsection{Разделение множественной нейронной селективности}
\label{sec:disentangling_selectivity}


Когда несколько поведенческих переменных коррелированы, видимая нейронная селективность
к одной переменной может возникать из-за её связи с другой \citep{mix,mix3}.
Для определения того, отражает ли перекрытие селективностей истинное совместное кодирование или
избыточность переменных, к парам поведенческих переменных применялась та же процедура оценки значимости на основе ВИ
(раздел~\ref{sec:statistical_assessment}). Если значимая ВИ между двумя поведенческими
переменными не обнаруживалась, нейроны, селективные к обеим, считались независимо информативными относительно каждой из них.

В случае, когда поведенческие переменные оказывались значимо связаны, селективность соответствующей нейронной активности далее анализировалась с помощью условной взаимной информации.
Пусть $A$ обозначает нейронную активность, а $X$ и $Y$ --- две поведенческие переменные. Условная
взаимная информация $I(A,X|Y)$ вычислялась в зависимости от типов $X$ и $Y$:

1. \textbf{$X$ и $Y$ непрерывные}:
\begin{equation*}
I(A,X|Y) = H(A,Y) + H(X,Y) - H(A,X,Y) - H(Y)
\end{equation*}
Энтропийные слагаемые оценивались методом GCMI с использованием разложения Холесского по ранг-нормализованным данным.

2. \textbf{$X$ непрерывная, $Y$ дискретная}:

Вычислялась $I(A,X)$ для каждого дискретного значения $Y$ по формуле:
\begin{equation*}
I(A,X) = H(A) + H(X) - H(A,X)
\end{equation*}
Полученные значения взвешивались эмпирическими вероятностями каждого значения $Y$.

3. \textbf{$Y$ непрерывная, $X$ дискретная}:
\begin{equation*}
I(A,X|Y) = I(A,X) - (I(A,Y) - I(A,Y|X))
\end{equation*}
Здесь $I(A,X)$ и $I(A,Y)$ вычислялись методом GCMI для переменных смешанного типа, а
$I(A,Y|X)$ --- аналогично случаю 2.

4. \textbf{$X$ и $Y$ дискретные}:
\begin{equation*}
I(A,X|Y) = H(A,Y) + H(X,Y) - H(A,X,Y) - H(Y)
\end{equation*}
$H(A,Y)$ вычислялась отдельно для каждого значения $Y$ с последующим усреднением; $H(X,Y)$ и
$H(Y)$ вычислялись по стандартным формулам дискретной энтропии, а $H(A,X,Y)$
оценивалась как сумма условных энтропий по всем комбинациям $X$, $Y$.
Затем вычислялась информация взаимодействия (II) \cite{McGill1954,ghassami2017},
количественно характеризующая избыточность или синергию в тройке $A$, $X$ и $Y$:
\begin{equation*}
II(A,X,Y) = I(A;X) - I(A;X|Y)
\end{equation*}
Данная величина симметрична относительно трёх переменных и усреднялась
по различным перестановкам их порядка. Ненулевое значение $|II|$ указывает на то, что
три попарные связи не являются независимыми: положительные значения возникают, когда
обусловливание уменьшает взаимную информацию (избыточность), а отрицательные ---
когда обусловливание увеличивает её (синергия) \cite{Bell2003}.

Следуя \cite{ghassami2017}, $|II|$ использовался как порог для выявления слабых
информационных путей в рамках направленного ациклического графа. Попарная связь
считалась <<слабой>>, если её взаимная информация была меньше $|II|$. При данном
ограничении не более одного из трёх попарных путей может быть слабым, что ограничивает
топологию взаимодействия двумя возможными конфигурациями (см. \fig{fig:graphs}).
\begin{figure}
\centering
\includegraphics[width=0.5\linewidth]{figs/cycles.png}
\caption{Два возможных ациклических
графа взаимодействия между тремя рассматриваемыми переменными при условии
слабой связи между $A$ и $Y$ (по \cite{ghassami2017}).
}
\label{fig:graphs}
\end{figure}
Процедура выполнялась следующим образом:
1. Вычислялась II для тройки $A$, $X$, $Y$.
2. Сравнивались MI($A,X$) и MI($A,Y$) с абсолютным значением II.
3. Если одна из связей являлась слабой, две другие считались сильными. В этом
случае одна поведенческая переменная фактически поглощала другую, а более сильная ВИ
интерпретировалась как содержательная (случай 1).
4. Если ни одна из связей не являлась слабой, переменные, вероятно, были избыточными, и для принятия решения о том, какую переменную сохранить, требовалась экспертная оценка (случай 2).

Например, если две поведенческие переменные $X$ и $Y$ коррелированы, процедура
определяла, какую из них нейрон преимущественно кодирует (случай 1), либо пара маркировалась как неоднозначная при данном наборе переменных (случай 2). В последнем случае обе переменные сохранялись для представления результатов, а для разрешения зависимости рекомендовалось включение дополнительных ковариат или экспериментальных манипуляций.

\subsubsection{Статистический анализ распределений задержек и селективностей}
\label{sec:statistical_analysis_delays}

Статистический анализ распределений задержек и количества селективных нейронов
выполнялся в GraphPad Prism 10.4.0 (GraphPad Software, США) с использованием
однофакторного дисперсионного анализа (ANOVA) с апостериорным тестом Тьюки.
Уровни значимости обозначены звёздочками:
$*$\,$p<0.05$, $**$\,$p<0.01$, $***$\,$p<0.001$, $****$\,$p<0.0001$;
незначимые пары опущены.
Столбцы отображают среднее $\pm$ SEM по комбинациям <<мышь--сессия>>.

\subsection{Анализ клеток места}
\label{sec:place_cells_analysis}

\subsubsection{Классический анализ клеток места}
\label{sec:classic_pc_identification}
Клетки места идентифицировались с помощью ранее описанной авторской процедуры в MATLAB
\cite{Plusnin2021}. Арена делилась на квадратные ячейки размером 8$\times$8 см. Для каждого нейрона
с числом кальциевых событий не менее $n=5$ строилась карта пространственной активности путём
деления сглаженного числа кальциевых событий на сглаженное время пребывания в каждой ячейке. Полученные карты сегментировались на различные пространственные области с помощью порогового метода на основе водораздела, после чего для каждой области вычислялось пространственное информационное содержание.
Для оценки значимости пространственной настройки генерировалось 1000 проб с перемешиванием путём
случайного сдвига следа кальциевой активности каждого нейрона относительно траектории животного. Пространственная информация для каждой пробы с перемешиванием пересчитывалась, и эмпирическое распределение использовалось для вычисления $p$-значения наблюдаемой
информации. Области, в которых пространственная информация превышала 95-й процентиль
распределения при перемешивании, рассматривались как кандидаты на поля места. Внутри каждой из них
поля места определялись как подобласти, в которых частота кальциевых событий превышала 50\%
от пиковой частоты разрядов в соответствующей области.


\subsubsection{Идентификация клеток места методом INTENSE}
\label{sec:intense_pc_identification}
Для оценки пространственной селективности в рамках INTENSE совместные пространственные
координаты животного $C={X,Y}$ рассматривались как поведенческая переменная. Все
информационно-теоретические меры вычислялись для объединённой пары координат,
включая энтропию $H(C) = H(X,Y)$ и взаимную информацию $I(A,C) = I(A; X,Y)$.
При процедуре перемешивания обе компоненты $C$ сдвигались совместно с использованием
одного и того же временного сдвига для сохранения координатной структуры.


\subsubsection{Сравнение популяций клеток места}
\label{sec:comparison_pc_populations}
Для сравнения популяций клеток места, идентифицированных методами INTENSE и классическим методом,
вычислялись метрики перекрытия по всем сессиям записи. Классический метод
идентифицировал клетки места по пространственному информационному содержанию с пороговым значением $z$-оценки
($z>3$) на основе пространственно разбитых частот разрядов, вычисленных по кальциевым событиям, обнаруженным с помощью
непрерывного вейвлет-преобразования \cite{Neugornet2021}. INTENSE идентифицировал
клетки места по взаимной информации между кальциевым сигналом флуоресценции и
пространственным положением (координаты $x$, $y$) со статистической значимостью, оценённой с помощью
двухэтапного пермутационного тестирования.

Перекрытие популяций количественно оценивалось с помощью трёх взаимодополняющих метрик. Первая
метрика --- пересечение относительно INTENSE --- представляла долю клеток места, идентифицированных INTENSE,
которые также были идентифицированы классическим методом, и вычислялась как $|A \cap B| /
|A|$, где $A$ --- клетки места по INTENSE, а $B$ --- клетки места по классическому методу. Вторая метрика --- пересечение относительно классического метода --- измеряла долю клеток места, идентифицированных
классическим методом, которые также были идентифицированы INTENSE, и вычислялась как $|A
\cap B| / |B|$. Третья метрика --- коэффициент Жаккара --- представляла симметричную
меру согласия и вычислялась как отношение пересечения к объединению: $|A \cap B| /
|A \cup B|$.

Для понимания влияния статистической достоверности на перекрытие популяций данные метрики оценивались
при различных порогах значимости от $p = 1.0$ до $p =
10^{-5}$. Для каждого порогового значения сначала применялся порог как к $p$-значениям INTENSE,
полученным из подобранного гамма-распределения с дополнительной массой в нуле, так и к классическим $p$-значениям, вычисленным
из $z$-оценок с помощью одностороннего нормального распределения. Затем в обоих методах сохранялись только нейроны с $p$-значениями ниже порога, и по отфильтрованной популяции вычислялись все три метрики
перекрытия. Дополнительно отслеживалась доля нейронов, сохранённых после фильтрации, для оценки строгости каждого порога.

Для анализа метки клеток места и статистические характеристики из всех 64
сессий записи были объединены, к объединённому набору данных применялись единые пороги, и
вычислялись глобальные метрики перекрытия.

\subsubsection{Сравнение пространственной селективности}
\label{sec:comparison_spatial_selectivity}

Проводилось сравнение карт пространственной активности, полученных методами INTENSE и классическим методом,
для оценки методологической конвергенции.

Сравнение карт пространственной активности требовало построения карт полей места
на основе как кальциевой флуоресценции, так и детектированных событий. Поведенческая
арена разделялась на сетку $7 \times 7$ пространственных ячеек --- разрешение, обеспечивающее баланс между пространственной
детализацией и достаточным числом отсчётов в ячейке при типичной продолжительности записи. Для каждой
пространственной ячейки вычислялась средняя нейронная активность в период пребывания животного в данном
местоположении, взвешенная по времени пребывания для компенсации неравномерного пространственного сэмплирования,
характерного для экспериментов со свободным поведением.

Для карт активности на основе кальциевого сигнала к следу флуоресценции предварительно применялся порог на уровне 80-го процентиля
для выделения периодов повышенной активности. Данный процентильный
порог адаптируется к динамическому диапазону каждого нейрона, фокусируясь на
надпороговых ответах, с наибольшей вероятностью соответствующих кальциевым
транзиентам, связанным с потенциалами действия. Пороговый сигнал затем усреднялся в пределах каждой пространственной ячейки за периоды пребывания. Карты активности на основе событий строились аналогично, но с использованием
дискретных моментов времени событий, обнаруженных с помощью вейвлетной детекции
(раздел~\ref{sec:calcium_event_detection}).

Для количественной оценки сходства карт полей места, полученных по кальциевому сигналу и по событиям, использовался
индекс структурного сходства (SSIM) --- перцептивно мотивированная метрика, изначально
разработанная для оценки качества изображений \cite{Wang2004ssim}. Обе карты активности
нормализовались к диапазону $[0, 1]$ перед сравнением. SSIM объединяет информацию
о яркости, контрасте и структурном сходстве, что делает эту метрику более робастной, чем
попиксельные сравнения, для оценки общего соответствия пространственных
карт. SSIM вычислялся с диапазоном данных 1.0, соответствующим нормализованным картам.

Дополнительно рассчитывались комплементарные метрики для всесторонней
оценки сходства полей места. Среднеквадратическая ошибка (MSE) количественно оценивала
среднюю квадратичную разность между нормализованными картами активности, обеспечивая
простую меру попиксельного отклонения. Пиковое отношение сигнал/шум
(PSNR) выражало эту разность относительно максимально возможного сигнала, предоставляя
масштабно-инвариантную метрику.

\subsubsection{Генерация поведенческих признаков}
\label{sec:synthetic_data_generation}

Для систематической оценки производительности INTENSE и сравнения с существующими методами
были сгенерированы синтетические нейронные наборы данных с известной эталонной связностью между нейронами и поведенческими признаками.

Были сгенерированы два типа поведенческих признаков для моделирования реалистичных экспериментальных
условий:

\textbf{Дискретные признаки:} Бинарные временные ряды, представляющие дискретные поведенческие
состояния (например, стойка на задних лапах, взаимодействие с объектом), генерировались с контролируемой временной
структурой. Для каждого признака создавались островки активности (единицы), перемежающиеся с
неактивными периодами (нулями), со следующими параметрами:
\begin{itemize}
\item Среднее число активных периодов: 10 на признак
\item Средняя продолжительность активного периода: 5 секунд (100 кадров при 20 Гц)
\item Экспоненциально распределённые промежутки между активными периодами
\item Продолжительность активных периодов из нормального распределения ($\sigma =
\text{длительность}/3$)
\end{itemize}

\textbf{Непрерывные признаки:} Непрерывные поведенческие переменные (например, аналоги скорости, направления
головы) генерировались с использованием дробного броуновского движения (fBM) со следующими параметрами:
\begin{itemize}
\item Показатель Хёрста $H = 0.3$ (субдиффузионное, антиперсистентное поведение)
\item Длина, соответствующая продолжительности эксперимента (1200 секунд)
\item Метод Дэвиса--Харта для вычислительно эффективной генерации
\item Каждый признак инициализировался уникальным случайным зерном для обеспечения независимости
\end{itemize}

\subsubsection{Генерация сигналов}
\label{sec:signal_generation}

\textbf{Связь <<признак--нейрон>>:} Каждый синтетический нейрон случайным образом назначался
ровно одному поведенческому признаку, что создавало разреженную эталонную матрицу связности.
Для непрерывных признаков выделялась область интереса (ROI), составляющая 15\% от
диапазона значений и центрированная на случайно выбранном процентиле, которая определяла <<рецептивное поле>> нейрона.

\textbf{Генерация событий с контролируемой надёжностью:} Нейронные события генерировались
с помощью двухрежимного неоднородного пуассоновского процесса \citep{Borst1999}:
\begin{itemize}
\item Базовая частота: $r_0 = 0.1$ Гц в неактивные периоды / вне ROI
\item Активная частота: $r_1 = 1.0$ Гц в активные периоды / внутри ROI
\item Вероятность пропуска: $p_{\text{skip}} \in [0, 1)$ для моделирования вариабельности ответов
\item События генерировались с частотой кадрирования записи (20 Гц)
\end{itemize}

Параметр вероятности пропуска $p_{\text{skip}}$ управляет надёжностью нейронного ответа
посредством случайного удаления активных периодов с вероятностью $p_{\text{skip}}$,
моделируя стохастическую природу нейронных ответов, наблюдаемую in vivo
\citep{Grijseels2021}.

\textbf{Кальциевая динамика:} Дискретные события разрядов свёртывались с
двойным экспоненциальным кальциевым ядром для генерации следов флуоресценции
\citep{Zhang2023}:
\begin{equation}
\label{eq:calcium_dynamics}
F(t) = \sum_i A_i \left(1 - \exp\!\left(-\frac{t - t_i}{\tau_r}\right)\right)
\exp\!\left(-\frac{t - t_i}{\tau_d}\right) \Theta(t - t_i) + \eta(t)
\end{equation}
где:
\begin{itemize}
\item $t_i$ --- моменты времени событий пуассоновского процесса
\item $A_i \sim \mathcal{U}(0.5, 2.0)$ --- амплитуды событий в единицах $\Delta F/F$
\item $\tau_r = 0.25$ секунды (постоянная времени нарастания)
\item $\tau_d = 2.0$ секунды (постоянная времени затухания), соответствующая кинетике GCaMP6s
\item $\Theta$ --- функция Хевисайда
\item $\eta(t) \sim \mathcal{N}(0, \sigma_{\text{noise}}^2)$ --- аддитивный гауссов шум
\end{itemize}

\textbf{Уровни шума:} Для оценки робастности варьировалось отношение сигнал/шум
(SNR), определяемое как отношение частот разрядов внутри и вне
целевых областей нейрона:

$\text{SNR} = \frac{r_1}{r_0}$

где:
- $r_0$ --- базовая частота разрядов вне области интереса (0.1 Гц);
- $r_1$ --- частота разрядов в пределах области интереса ($r_1 = r_0 \times \text{SNR}$).

Например, SNR = 8 означает, что нейрон разряжается с частотой 0.8 Гц, когда поведенческая переменная
находится в пределах его рецептивного поля, по сравнению с базовой частотой 0.1 Гц. Данное определение
позволяет систематически варьировать силу нейронной селективности от слабой
(SNR = 2, пиковая частота 0.2 Гц) до очень сильной (SNR = 64, пиковая частота 6.4 Гц).

Использовались значения SNR из ряда $\{2, 4, 8, 16, 32, 64\}$, охватывающего диапазон
от сильно зашумлённых до практически бесшумных условий.

\subsubsection{Состав валидационного набора данных}
\label{sec:validation_dataset}

Каждый синтетический набор данных включал:
\begin{itemize}
\item 500 нейронов (250 с дискретной селективностью, 250 с непрерывной селективностью)
\item 20 поведенческих признаков (10 дискретных, 10 непрерывных)
\item Эталонная связность: 5\% (500 истинных связей из 10\,000 возможных пар)
\item Длительность записи: 1200 секунд (24\,000 кадров при 20 Гц)
\item Сетка параметров: 6 значений SNR $\times$ 9 значений $p_{\text{skip}}$ = 54 условия
\end{itemize}

\subsubsection{Метрики оценки производительности}
\label{sec:performance_metrics}

Для количественной оценки качества детекции вычислялись стандартные метрики бинарной классификации путём сравнения обнаруженных значимых пар <<нейрон--признак>> с эталонными данными:

\textbf{Точность (Precision):} Доля обнаруженных связей, являющихся истинно положительными:
\begin{equation}
\label{eq:precision}
\text{Precision} = \frac{\text{TP}}{\text{TP} + \text{FP}}
\end{equation}

\textbf{Полнота (Recall):} Доля истинных связей, которые были обнаружены:
\begin{equation}
\label{eq:recall}
\text{Recall} = \frac{\text{TP}}{\text{TP} + \text{FN}}
\end{equation}

\textbf{F1-мера:} Гармоническое среднее точности и полноты:
\begin{equation}
\label{eq:f1score}
\text{F1} = 2 \cdot \frac{\text{Precision} \times \text{Recall}}{\text{Precision} +
\text{Recall}}
\end{equation}

где TP --- истинно положительные, FP --- ложноположительные, FN --- ложноотрицательные результаты.

\textbf{Случайная базовая линия:} При эталонной связности 5\% случайный классификатор
достигал бы:
\begin{equation}
\label{eq:random_baseline}
\text{Precision}_{\text{random}} = 0.05
\end{equation}

\subsubsection{Методы сравнения}
\label{sec:comparison_methods}

INTENSE сравнивался с тремя категориями подходов к оценке связи <<нейрон--поведение>>:

\textbf{1. Простые критериальные методы (без перемешивания):}
\begin{itemize}
\item \textbf{Для непрерывных признаков:} Коэффициент корреляции Пирсона со значимостью, определяемой по $t$-критерию Стьюдента. Для выборки размера $n$ тестовая статистика:
\begin{equation}
\label{eq:pearson_test}
t = r\sqrt{\frac{n-2}{1-r^2}}
\end{equation}
имеет $t$-распределение с $n-2$ степенями свободы при нулевой гипотезе о нулевой корреляции \citep{Fisher1915}.

\item \textbf{Для дискретных признаков:} $t$-критерий Уэлча для сравнения средней кальциевой активности
между активными и неактивными поведенческими периодами. Данный критерий учитывает неравенство
дисперсий между группами и использует аппроксимацию Уэлча--Саттертуэйта для степеней
свободы \citep{Welch1947}:
\begin{equation}
\label{eq:welch_df}
\nu = \frac{\left(\frac{s_1^2}{n_1} +
\frac{s_2^2}{n_2}\right)^2}{\frac{s_1^4}{n_1^2(n_1-1)} + \frac{s_2^4}{n_2^2(n_2-1)}}
\end{equation}
\end{itemize}

\textbf{2. Методы на основе ВИ (без перемешивания):}
\begin{itemize}
\item \textbf{Для дискретных признаков:} использовался подставной (plug-in) оценщик с
адаптивным разбиением. Число интервалов выбиралось по расширенному правилу Стёрджеса:
$n_{\text{bins}} = \lceil \log_2(n) \rceil + 1$ \citep{Sturges1926}, с
квантильным разбиением для обеспечения равного числа наблюдений в каждом интервале.

Статистическая значимость оценивалась с помощью критерия отношения правдоподобия ($G$-критерий
независимости) \citep{Sokal1995}:

\begin{equation}
\label{eq:gtest}
G = 2 \sum_{i,j} O_{ij} \ln\left(\frac{O_{ij}}{E_{ij}}\right)
\end{equation}

где $O_{ij}$ --- наблюдаемые частоты, а $E_{ij}$ --- ожидаемые частоты при
независимости. Тестовая статистика имеет распределение $\chi^2$ с $(r-1)(c-1)$
степенями свободы для $r$ строк и $c$ столбцов таблицы сопряжённости.

\item \textbf{Для непрерывных признаков:} ВИ преобразовывалась в эффективную корреляцию
по соотношению \citep{Cover2006}:
\begin{equation}
\label{eq:mi_to_corr}
\rho_{\text{eff}} = \text{sign}(\rho_{\text{sample}}) \cdot \sqrt{1 - \exp(-2 \cdot
\text{MI})}
\end{equation}
Это позволяло вычислить $p$-значение с помощью стандартного $t$-критерия для корреляции.
\end{itemize}

\textbf{3. Конвейеры типа INTENSE:}
\begin{itemize}
\item \textbf{На основе корреляции:} Конвейер INTENSE с использованием векторизованной корреляции Пирсона
в качестве метрики сходства. Сохраняя полную статистическую строгость
двухэтапной процедуры, данный тест оценивает линейную зависимость между сигналами.

\item \textbf{На основе средних:} Конвейер INTENSE с использованием разности средних в качестве
метрики для дискретных признаков. Тестовая статистика представляет собой разность средней кальциевой
активности между поведенческими состояниями.

\item \textbf{На основе ВИ (INTENSE):} Полный конвейер INTENSE с взаимной информацией на основе гауссовой копулы
(GCMI) \citep{gcmi}. GCMI преобразует переменные в представление гауссовой копулы, где ВИ имеет замкнутую форму (ур.~\ref{eq:gcmi_continuous}).
Данный подход устойчив к выбросам и не требует параметров разбиения.
\end{itemize}

Все конвейеры типа INTENSE использовали двухэтапную процедуру тестирования, описанную в
разделе~\ref{sec:statistical_assessment}, с идентичными параметрами (100 перемешиваний для скрининга,
10\,000 перемешиваний для валидации, FWER = 0.01).

\subsubsection{Протокол эксперимента}
\label{sec:experimental_protocol}

Для каждой комбинации параметров (SNR, $p_{\text{skip}}$) выполнялось 5 независимых
случайных инициализаций, а полученные значения точности, полноты и F1-меры усреднялись.
Данное усреднение снижает влияние случайных флуктуаций и обеспечивает более стабильные
оценки производительности. Всесторонняя оценка охватывала:
\begin{itemize}
\item 6 значений SNR: $\{2, 4, 8, 16, 32, 64\}$
\item 9 значений $p_{\text{skip}}$: $\{0.0, 0.1, 0.2, 0.3, 0.4, 0.5, 0.6, 0.7, 0.8\}$
\item Всего: $6 \times 9 \times 5 = 270$ синтетических наборов данных на тип признака
\end{itemize}

Ко всем методам применялось понижение частоты дискретизации в 5 раз для снижения вычислительной нагрузки
при сохранении достаточного временного разрешения для обнаружения ассоциаций <<нейрон--признак>>.
